\section{Introduction}
\label{sec:introduction}

In the financial year ending 2025, 31 per cent of UK children were in relative low income after housing costs, and 13 per cent were in deep material poverty -- lacking four or more of the items the Department for Work and Pensions treats as essential to a child's standard of living \citep{dwp2026,jrf2026}. Since low household income does not always mean a child lacks specific essentials -- a warm coat, a proper meal, a school trip -- these two measures diverge. Income poverty is more than twice as common as deep material deprivation. That distinction matters for policy as raising household income does not guarantee a proportional fall in material hardship.

Child poverty has also become a more explicit UK Government priority in recent years. Its Child Poverty Strategy, published December 2025, centres on scrapping the two-child Universal Credit limit from April 2026. This has been described as freeing around 1.6 million children from poverty \citep{jrf2026}.

Social security is largely reserved to the UK Government, but Scotland has gone furthest among the devolved administrations in topping it up, through the Scottish Child Payment (SCP), described on launch by \citet{scottishgovernment2021} as "the most ambitious anti-poverty measure currently being undertaken anywhere in the UK." SCP is a UK-first, a benefit unique to Scotland with no equivalent elsewhere in Great Britain \citep{scottishgovernment2021,andersen2025}, and the centrepiece of the Scottish Government's response to statutory targets in the \citet{childpovertyact2017}, which commit Scotland to cutting relative child poverty below 10 per cent, and combined low income and material deprivation below 5 per cent, by 2030. Launched for under-6s in February 2021 and extended to all under-16s from November 2022 (Section~\ref{sec:data} gives the exact dates and rates), its weekly rate has more than doubled since launch, to \pounds28.20 from April 2026 \citep{scottishgovernment2026}. It sits alongside four related Best Start payments delivered by Social Security Scotland; together these are worth around \pounds25,000 to an eligible family by a child's 16th birthday, against under \pounds2,000 in England and Wales, where equivalent support ends at age four \citep{scottishgovernment2026}. Scottish Government modelling credits SCP alone with keeping around 40,000 children out of relative poverty in 2025--26 \citep{scottishgovernment2025}. That modelling, like \citet{andersen2025}'s comparison discussed below, relies on a counterfactual rather than the kind of quasi-experimental design used in this paper. It is a live test of whether a devolved government can reduce child poverty and material hardship independently of UK-wide policy.

In July 2026, the Welsh Government's own expert group met for the first time to design Cynnal, a proposed \pounds10-a-week child payment pilot for up to 15,000 children in Universal Credit households \citep{welshgovernment2026}. Whether such a transfer meaningfully reduces material deprivation, not just income poverty, is therefore a live design question facing at least two devolved nations, not one.

The closest existing evidence is \citet{andersen2025}, reviewed in full in Section~\ref{sec:literature-review}. \citeauthor{andersen2025} find that SCP reduced material deprivation, but from a single, linear OLS estimate. This poses the question as to whether that finding survives re-estimation with a less restrictive, machine-learning estimator, and whether it reflects genuine change across specific hardships -- a warm coat, a proper meal, a school trip -- rather than an aggregate masking null effects once each is tested on its own. 

This paper answers both. Using nine years of Family Resources Survey/HBAI data, it first replicates \citet{andersen2025}'s OLS specification, then re-estimates the same effect with a cross-fitted, doubly robust machine-learning estimator, and decomposes the result across the ten items using three independent estimation strategies. The headline finding is that SCP produced a real, but more modest, reduction in material deprivation than \citeauthor{andersen2025}'s OLS estimate alone implies, with no single hardship item showing a robust effect once correctly adjusted for. Section~\ref{sec:literature-review} reviews the relevant literature; Section~\ref{sec:data} describes the data and outcome measures; Section~\ref{sec:methodology} sets out the empirical strategy; Section~\ref{sec:results} presents the results; and the discussion and conclusion that follow assess what these findings mean for the design of Cynnal and any future UK child payment.